\section{DevOps and DevSecOps}
\label{sec:DevSecOps}

DevOps emerged as a software engineering methodology that promotes close collaboration between software development and IT operations through automation, continuous integration, continuous delivery, and continuous monitoring. Rather than treating development and operations as independent activities, DevOps integrates them into a continuous workflow that accelerates software delivery while improving consistent deployment and operational efficiency \cite{Bass2015DevOps},\cite{Humble2010ContinuousDelivery}.

The cloud computing paradigm has further accelerated the adoption of DevOps. Cloud platforms provide on-demand infrastructure, infrastructure-as-code (IaC), and automated deployment services that naturally is compatible with DevOps practices. These capabilities enable organizations to provision infrastructure rapidly, automate software deployment, and monitor applications across distributed environments. Consequently, DevOps has become one of the predominant software delivery methodologies for cloud-native and hybrid cloud applications \cite{Bass2015DevOps}, \cite{Liu2011NISTRA}.

Despite its advantages, DevOps primarily emphasizes development velocity, automation, and operational efficiency, while security activities have traditionally remained separate from the continuous development pipeline. Security assessments are frequently performed after software implementation or before deployment, increasing development costs and delaying software releases. As organizations deploy applications across cloud and hybrid cloud infrastructures, this separation becomes more problematic because security policies, infrastructure configurations, software dependencies, and cloud resources continuously evolve throughout the software operational lifecycle \cite{Rajapakse2021DevSecOpsSLR}, \cite{Myrbakken2017DevSecOps}.

To address these limitations, DevSecOps extends DevOps by integrating security practices into every stage of the software development and operational lifecycle. Rather than treating security as a separate verification phase, DevSecOps advocates continuous security through automated security testing, vulnerability assessment, compliance verification, infrastructure security, and continuous monitoring. The NIST Secure Software Development Framework (SSDF) similarly recommends incorporating security practices throughout software development to reduce vulnerabilities and improve software integrity \cite{NISTSSDF2022}. Recent literature reviews further identify continuous security integration, automation, and security orchestration as the defining characteristics of DevSecOps while highlighting that these remain active areas of research, particularly for complex cloud environments \cite{Rajapakse2021DevSecOpsSLR}, \cite{Zhao2024DevSecOpsMLR}.

Although DevSecOps significantly strengthens software security, existing research primarily concentrates on integrating security into software development workflows. Comparatively less attention has been devoted to systematically coordinating software development, infrastructure management, operational activities, and continuous security within heterogeneous hybrid cloud environments. This observation motivates the SysSecOps operational process proposed in this research, which extends existing DevSecOps practices by integrating system administration and infrastructure operations into a unified operational workflow for hybrid cloud environments.